% =============================================================================
% s74_branches.tex --- Pre-written replacement text for main.tex §7.4 (E5
% paragraph), branched on the v6 mask-weighted-L1 re-eval VERDICT.
%
% USAGE
% -----
% After `project/analyze_e5_perclass.py` finishes on the v6 (mask-L1,
% lambda_mask=3.0, 4x lesion weight) checkpoint, it prints a VERDICT line:
%
%   "mask-L1 HELPS melanoma"      -> use BRANCH A below
%   "mask-L1 does NOT help melanoma" -> use BRANCH B below
%
% Each branch is a drop-in replacement for main.tex lines 295-308, i.e. the
% \paragraph{E5 --- ...} block (everything from "\paragraph{E5" through the
% sentence ending "...left to future work (\cref{sec:discussion})."
% inclusive). Line 294 (\subsection{...}\label{sec:exp-salvage}) and the E6
% paragraph (310-315) are UNCHANGED in both branches.
%
% PLACEHOLDERS (fill in after running analyze_e5_perclass.py on v6):
%   SALV_MEL_V6  -- v6 melanoma salvage rate, e.g. "5.2\%" -> new value
%   DMG_MEL_V6   -- v6 melanoma damage rate (correctly-classified flipped wrong)
%   NET_V6       -- v6 net malignant case count, signed (e.g. "-81" -> "-40" or "+12")
%   SALV_MEL_V6_FRAC -- raw fraction form, e.g. "4/77" -> "x/77" (denominator may
%                        differ slightly if the v6 probe re-splits; check n)
%   DMG_MEL_V6_FRAC  -- raw fraction form, e.g. "85/274" -> "y/274"
%   NET_V6_ABS   -- absolute value of NET_V6 for prose ("a net loss of N" /
%                   "a net gain of N"), used only in Branch A success wording
%
% Frozen v5 numbers (DO NOT CHANGE, used for before/after comparison in both
% branches): benign salvage 75.6% (1809/2392), benign damage 0.6% (50/8138),
% melanoma salvage 5.2% (4/77), melanoma damage 31.0% (85/274), net -81,
% aggregate salvage rate 0.737, n=3627 per severity.
% =============================================================================


% =============================================================================
% BRANCH A --- mask-L1 HELPS melanoma (net rises from -81 toward / past zero)
% =============================================================================
% Use when VERDICT == "mask-L1 HELPS melanoma", i.e. NET_V6 > -81 (net loss
% shrinks) and ideally NET_V6 >= 0 (net gain). Two sub-variants below depending
% on the sign of NET_V6 --- pick ONE and delete the other before pasting into
% main.tex.
%
% -----------------------------------------------------------------------
% Variant A1 --- NET_V6 turns POSITIVE (net gain for malignant cases)
% -----------------------------------------------------------------------
% Use this wording if NET_V6 >= 0.

\paragraph{E5 --- mask-weighted reconstruction narrows the malignant salvage gap.}
On the moderate band the aggregate salvage rate --- the fraction of degraded-then-misclassified
images that enhancement restores to a correct prediction --- is $0.737$, nominally clearing the
$>55\%$ target. As in the feature-DP model, this aggregate is misleading on its own and we do
\emph{not} report it as a headline capability: a per-class split (independent EfficientNet-B3
probe, $n=3627$ per severity) is required to see what the aggregate hides. Under the
feature-DP enhancer, melanoma salvage was just $5.2\%$ ($4/77$) against $31.0\%$ damage
($85/274$ correctly-classified), a net loss of $81$ malignant cases, while benign lesions
salvaged $75.6\%$ ($1809/2392$) at $0.6\%$ damage ($50/8138$). With \emph{mask-weighted}
reconstruction ($\lambda_{\text{mask}}=3.0$, $4\times$ lesion-region weight, which forbids the
enhancer from smoothing over the lesion itself), melanoma salvage rises to
$\mathrm{SALV\_MEL\_V6}$ ($\mathrm{SALV\_MEL\_V6\_FRAC}$) while damage falls to
$\mathrm{DMG\_MEL\_V6}$ ($\mathrm{DMG\_MEL\_V6\_FRAC}$), turning the net malignant outcome from a
loss of $81$ cases into a \textbf{net gain of $\mathrm{NET\_V6\_ABS}$ cases}
($\mathrm{NET\_V6}$). This is consistent with the design hypothesis: protecting the lesion
region from over-smoothing is what lets enhancement help the class that matters most. We still
do \emph{not} promote the aggregate $0.737$ figure, and the benign-dominated shape of the
aggregate salvage rate is unchanged --- mask-weighted reconstruction shifts \emph{where} the
salvage comes from, not how the headline number should be read. The residual dangerous flips on
the moderate band (\cref{fig:dflip}) and the severe-band degradation below still motivate the
query-for-retake channel (\cref{sec:discussion}) as the primary safety mechanism; mask-weighted
reconstruction is a complementary mitigation for the band where enhancement \emph{is} the
selected action, not a replacement for the gate.

% -----------------------------------------------------------------------
% Variant A2 --- NET_V6 improves but stays NEGATIVE (net loss shrinks, doesn't
% flip sign)
% -----------------------------------------------------------------------
% Use this wording if NET_V6 < 0 but |NET_V6| < 81 (i.e. less negative than v5).

\paragraph{E5 --- mask-weighted reconstruction reduces but does not eliminate the malignant
salvage deficit.}
On the moderate band the aggregate salvage rate --- the fraction of degraded-then-misclassified
images that enhancement restores to a correct prediction --- is $0.737$, nominally clearing the
$>55\%$ target. As in the feature-DP model, this aggregate is misleading on its own and we do
\emph{not} report it as a headline capability: a per-class split (independent EfficientNet-B3
probe, $n=3627$ per severity) is required to see what the aggregate hides. Under the
feature-DP enhancer, melanoma salvage was just $5.2\%$ ($4/77$) against $31.0\%$ damage
($85/274$ correctly-classified), a net loss of $81$ malignant cases, while benign lesions
salvaged $75.6\%$ ($1809/2392$) at $0.6\%$ damage ($50/8138$). With \emph{mask-weighted}
reconstruction ($\lambda_{\text{mask}}=3.0$, $4\times$ lesion-region weight, which forbids the
enhancer from smoothing over the lesion itself), melanoma salvage rises to
$\mathrm{SALV\_MEL\_V6}$ ($\mathrm{SALV\_MEL\_V6\_FRAC}$) and damage falls to
$\mathrm{DMG\_MEL\_V6}$ ($\mathrm{DMG\_MEL\_V6\_FRAC}$), reducing the net malignant loss from
$81$ to $\mathrm{NET\_V6\_ABS}$ cases ($\mathrm{NET\_V6}$). Mask-weighted reconstruction thus
\emph{reduces but does not eliminate} the malignant-salvage deficit: protecting the lesion
region from over-smoothing measurably helps, but enhancement still does net harm to melanoma
cases on the moderate band. We do not promote the aggregate $0.737$ figure, which remains
benign-dominated. We read the residual deficit, together with the dangerous flips on the
moderate band (\cref{fig:dflip}) and the severe-band degradation below, as confirmation that a
query-for-retake gate (\cref{sec:discussion}) --- not loss reweighting alone --- is the
necessary safety mechanism for the malignant class; mask-weighted reconstruction is best framed
as a partial, complementary mitigation rather than a fix.


% =============================================================================
% BRANCH B --- mask-L1 does NOT help melanoma (net flat or worse than -81)
% =============================================================================
% Use when VERDICT == "mask-L1 does NOT help melanoma", i.e. NET_V6 <= -81
% (net loss unchanged or larger). This is essentially the current (v5) framing,
% reinforced with the v6 negative result as evidence that loss-reweighting
% alone cannot close the gap.

\paragraph{E5 --- salvage rate is benign-dominated, and mask-weighted reconstruction does not
close the gap.}
On the moderate band the aggregate salvage rate --- the fraction of degraded-then-misclassified
images that enhancement restores to a correct prediction --- is $0.737$, nominally clearing the
$>55\%$ target. This aggregate is misleading, and we deliberately do \emph{not} report it as a
headline capability: a per-class split (independent EfficientNet-B3 probe, $n=3627$ per severity)
shows it is almost entirely benign false-positive correction. On benign lesions enhancement salvages
$75.6\%$ ($1809/2392$) while damaging only $0.6\%$ ($50/8138$ correctly-classified); on melanoma it
salvages just $5.2\%$ ($4/77$) while damaging $31.0\%$ ($85/274$ correctly-classified) --- a net loss
of $81$ malignant cases. Enhancement thus helps the benign-leaning errors that least need help and
harms the malignant cases that matter most, mirroring the residual dangerous flips above
(\cref{fig:dflip}) and the severe-band degradation below. We tested whether this gap is an
artefact of the reconstruction loss: a mask-weighted variant ($\lambda_{\text{mask}}=3.0$,
$4\times$ lesion-region weight, forbidding the enhancer from smoothing over the lesion) yields
melanoma salvage of $\mathrm{SALV\_MEL\_V6}$ ($\mathrm{SALV\_MEL\_V6\_FRAC}$) and damage of
$\mathrm{DMG\_MEL\_V6}$ ($\mathrm{DMG\_MEL\_V6\_FRAC}$), a net of $\mathrm{NET\_V6}$ malignant
cases --- statistically indistinguishable from (or worse than) the feature-DP result. Reweighting
the reconstruction loss toward the lesion is therefore \emph{not sufficient} to close the
malignant salvage gap: the failure mode is not "the enhancer smooths away the lesion" but
something the per-pixel reconstruction objective cannot reach (e.g.\ the diagnostic decision
boundary depends on texture statistics that mask-weighted L1 does not target). We read E5 as the
strongest evidence that \emph{blanket} enhancement is unsafe and that a query-for-retake gate is
required --- not as a salvage benchmark to maximise. We report this negative result rather than
omit it; closing the malignant salvage gap by other means (e.g.\ diagnosis-aware or
adversarial objectives) is left to future work (\cref{sec:discussion}).
